% bismi ALLAH ARRAHMAN ARRAHIM










































\documentclass[11pt,a4paper]{article}

% -------------------------------------------------
% Encodage, langue, typographie
% -------------------------------------------------
\usepackage[french]{babel}
\usepackage{fontspec}
\usepackage{microtype}
\usepackage[a4paper,margin=2.5cm]{geometry}

% -------------------------------------------------
% Mathématiques
% -------------------------------------------------
\usepackage{amsmath,amssymb,amsthm,mathtools,bm}
\usepackage{enumitem}

% -------------------------------------------------
% Liens
% -------------------------------------------------
\usepackage[hidelinks]{hyperref}

% -------------------------------------------------
% Police
% -------------------------------------------------
\setmainfont{TeX Gyre Pagella}

% -------------------------------------------------
% Environnements
% -------------------------------------------------
\newtheorem{definition}{Définition}[section]
\newtheorem{assumption}{Hypothèse}[section]
\newtheorem{proposition}{Proposition}[section]
\newtheorem{theorem}{Théorème}[section]
\newtheorem{corollary}{Corollaire}[section]
\newtheorem{remark}{Remarque}[section]

% -------------------------------------------------
% Opérateurs
% -------------------------------------------------
\DeclareMathOperator{\prox}{prox}
\DeclareMathOperator{\argmin}{argmin}
\DeclareMathOperator{\supp}{supp}
\DeclareMathOperator{\dist}{dist}
\DeclareMathOperator{\KL}{KL}
\DeclareMathOperator{\Id}{Id}
\DeclareMathOperator{\diag}{diag}
\DeclareMathOperator{\softmax}{softmax}

% -------------------------------------------------
% Raccourcis
% -------------------------------------------------
\newcommand{\X}{\mathcal X}
\newcommand{\C}{\mathcal C}
\newcommand{\D}{\mathcal D}
\newcommand{\E}{\mathcal E}
\newcommand{\Z}{\mathcal Z}
\newcommand{\Hh}{\mathcal H}
\newcommand{\I}{\mathcal I}
\newcommand{\T}{\mathcal T}
\newcommand{\W}{\mathcal W}
\newcommand{\N}{\mathcal N}
\newcommand{\R}{\mathbb R}
\newcommand{\Pcal}{\mathcal P}
\newcommand{\Qcal}{\mathcal Q}
\newcommand{\Gcal}{\mathcal G}
\newcommand{\Mcal}{\mathcal M}
\newcommand{\1}{\mathbf 1}

\title{PCH-ESM\\
\large Purification Contextuelle Hiérarchique et Espaces Sémantiques Modulaires}
\author{}
\date{}

\begin{document}

\maketitle

\begin{abstract}
Ce document formalise une architecture hiérarchique pour modèles séquentiels et systèmes de langage, fondée sur trois piliers :
(i) la purification récursive d'une représentation dans un contexte donné,
(ii) la décomposition hiérarchique d'un espace sémantique global en sous-espaces spécialisés,
et (iii) une fonction de coordination \(G\) qui transfère la structure de trajectoire d'un modèle généraliste vers un ensemble de modèles spécialisés.
La théorie distingue explicitement la représentation, la stabilité contextuelle, la hiérarchie, la spécialisation, la coordination inter-espaces et le rôle de modèle du monde.
\end{abstract}

\tableofcontents

% =================================================
\section{Cadre général}
% =================================================

L'idée directrice de PCH-ESM est qu'une bonne représentation n'est pas seulement une représentation prédictive, mais une représentation stable \emph{dans un contexte}.
Cette stabilité n'est pas absolue : elle dépend de la tâche, du domaine, du mode d'exécution, des contraintes de calcul et de la structure locale de coopération entre sous-modèles.

On distingue donc trois niveaux :

\begin{enumerate}[leftmargin=2em]
    \item un niveau de purification hiérarchique de base, noté PH ;
    \item un niveau contextuel, noté PCH ;
    \item un niveau architectural distribué, noté Multi-ES.
\end{enumerate}

L'architecture complète combine ensuite une couche de coordination \(G\), puis éventuellement une couche de modèle du monde Multi-PH.

% =================================================
\section{PH : purification hiérarchique}
% =================================================

Cette section résume la base PH déjà présente dans le document EEG.

\begin{definition}[Opérateur de purification hiérarchique]
Soit \(u\) un objet séquentiel ou latent. On définit un opérateur de purification \(\Pcal_\alpha\), de paramètre \(\alpha\in[0,1]\), qui agit récursivement via deux passes directionnelles :
\[
\mathcal F_\alpha(u)_i
=
\alpha\,\mathcal F_\alpha(u)_{i-1}
+
(1-\alpha)u_i,
\]
\[
\mathcal B_\alpha(u)_i
=
\alpha\,\mathcal B_\alpha(u)_{i+1}
+
(1-\alpha)u_i,
\]
et la purification symétrique :
\[
\Pcal_\alpha(u)
=
\frac12\Bigl(\mathcal F_\alpha(u)+\mathcal B_\alpha(u)\Bigr).
\]
\end{definition}

\begin{definition}[Hiérarchie de noyaux]
On définit :
\[
N^{(0)} = W,
\qquad
N^{(r+1)} = \Pcal_\alpha\!\left(N^{(r)}\right),
\quad r\ge 0.
\]
Le noyau terminal est
\[
N^{(\star)} = N^{(r^\star)}
\]
lorsque la suite converge ou lorsqu'une profondeur maximale est atteinte.
\end{definition}

\begin{definition}[Résidu]
Le résidu associé au noyau terminal est
\[
R^{(0)} = W - N^{(\star)}.
\]
Il peut être purifié à son tour, possiblement sous guidage du noyau terminal.
\end{definition}

\begin{proposition}[Attention standard comme cas plat]
L'attention standard peut être vue comme un cas plat de purification hiérarchique lorsque la hiérarchie a profondeur \(1\), que la sélection est bilinéaire, et que la fusion est une moyenne pondérée des valeurs.
Dans ce cas, l'opérateur \(QK^\top\) n'est pas la théorie fondamentale mais une implémentation particulière d'une sélection contextuelle plate.
\end{proposition}

% =================================================
\section{De PH vers PCH : purification contextuelle hiérarchique}
% =================================================

Cette section généralise PH en introduisant explicitement un contexte.

\subsection{Espace des contextes}

\begin{definition}[Espace des contextes]
On note \(\C\) l'espace des contextes.
Un contexte est un élément
\[
c\in\C.
\]
On peut décomposer un contexte sous la forme
\[
c=(d,\tau,m,\omega,\xi),
\]
où :
\begin{itemize}[leftmargin=2em]
    \item \(d\) désigne le domaine ou la spécialité ;
    \item \(\tau\) désigne la tâche ou l'objectif ;
    \item \(m\) désigne le mode d'exécution ou de coopération ;
    \item \(\omega\) désigne des contraintes de calcul, mémoire ou temps ;
    \item \(\xi\) désigne des paramètres additionnels propres à l'application.
\end{itemize}
\end{definition}

\begin{remark}
Le contexte n'est pas une simple étiquette. Il encode la manière dont l'objet doit être interprété, purifié et utilisé.
Deux contextes distincts peuvent donc induire deux stabilités différentes pour une même représentation.
\end{remark}

\subsection{Stabilité contextuelle}

\begin{definition}[Fonctionnelle de stabilité contextuelle]
Soit \(\Z\) un espace latent.
On définit une fonctionnelle
\[
J:\Z\times\C\to\R\cup\{+\infty\},
\]
où \(J(z;c)\) mesure le coût d'instabilité de la représentation \(z\) dans le contexte \(c\).
On suppose que, pour chaque \(c\in\C\), la fonction \(z\mapsto J(z;c)\) est propre, semi-continue inférieurement, et au moins coercive sur le sous-espace considéré.
\end{definition}

\begin{definition}[Stabilité contextuelle]
Une représentation \(z^\star\in\Z\) est dite \emph{stable dans le contexte \(c\)} si
\[
z^\star\in\argmin_{z\in\Z} J(z;c).
\]
Si l'ensemble des minimisateurs est singleton, on parle de \emph{noyau contextuel unique}.
Sinon, on parle d'un \emph{ensemble de noyaux contextuels}.
\end{definition}

\begin{remark}
La stabilité n'est donc plus globale ni absolue : elle est relative au contexte \(c\).
C'est précisément ce qui distingue PCH de PH.
\end{remark}

\subsection{Opérateur de purification contextuelle}

\begin{definition}[Opérateur de purification contextuelle]
Pour un pas \(\lambda>0\), on définit
\[
\Pcal_{\lambda}(\,\cdot\,;c):\Z\to\Z,
\qquad
\Pcal_{\lambda}(z;c)
=
\argmin_{u\in\Z}
\left\{
J(u;c)
+
\frac{1}{2\lambda}\|u-z\|^2
\right\}.
\]
\end{definition}

\begin{remark}
Lorsque \(J(\cdot;c)\) est convexe, \(\Pcal_{\lambda}(\cdot;c)\) est l'opérateur proximal associé à la stabilité contextuelle.
Lorsque \(J\) est non convexe, la définition reste formelle mais nécessite un choix de branche ou une approximation numérique.
\end{remark}

\subsection{Trajectoires de purification contextuelle}

\begin{definition}[Trajectoire de purification contextuelle]
Soit une entrée \(x\in\X\), un encodeur \(E_\theta:\X\to\Z\), et un contexte \(c\in\C\).
La trajectoire de purification contextuelle associée à \((x,c)\) est la suite
\[
z^{(0)}(x,c)=E_\theta(x),
\qquad
z^{(r+1)}(x,c)=\Pcal_{\lambda}\!\bigl(z^{(r)}(x,c);c\bigr),
\quad r\ge 0.
\]
On note cette trajectoire
\[
\T(x,c)=\bigl(z^{(0)}(x,c),z^{(1)}(x,c),z^{(2)}(x,c),\dots\bigr).
\]
\end{definition}

\begin{definition}[Noyau contextuel]
Si la suite \(\{z^{(r)}(x,c)\}_{r\ge 0}\) converge, on définit le noyau contextuel par
\[
N^\star(x,c)=\lim_{r\to\infty}z^{(r)}(x,c).
\]
Si la convergence numérique n'est pas garantie, on prend une approximation au rang \(r^\star\) :
\[
N^\star(x,c)\approx z^{(r^\star)}(x,c).
\]
\end{definition}

\begin{definition}[Résidu contextuel]
Le résidu contextuel associé à \((x,c)\) est
\[
R^{(0)}(x,c)=z^{(0)}(x,c)-N^\star(x,c).
\]
Il représente la partie de la représentation initiale qui n'est pas absorbée par la stabilité contextuelle principale.
\end{definition}

\begin{theorem}[PH est un cas limite de PCH]
Supposons que le contexte soit trivial au sens suivant :
\[
\mathcal M_c = \Z
\quad\text{pour tout } c\in\C,
\]
et que le décodeur soit contextuellement inactif, c'est-à-dire
\[
D_\psi(z,c)=D_\psi(z)
\quad\text{pour tout } z\in\Z,\; c\in\C.
\]
Alors la pénalité contextuelle s'annule identiquement :
\[
\ell_{\mathrm{ctx}}(z;c)=0,
\]
et la fonctionnelle fondamentale devient indépendante du contexte :
\[
J(z;x,c)=J(z;x).
\]
Par conséquent, l'opérateur \(\Pcal_{\lambda}(z;x,c)\) se réduit à l'opérateur PH classique.
\end{theorem}

\begin{proof}
Si \(\mathcal M_c=\Z\), alors
\[
\dist(z,\mathcal M_c)=0
\]
pour tout \(z\in\Z\), donc \(\ell_{\mathrm{ctx}}(z;c)=0\).
Si en outre \(D_\psi(z,c)=D_\psi(z)\), la dépendance à \(c\) disparaît complètement de \(J\).
L'opérateur proximal associé devient alors exactement celui de PH.
\end{proof}

\begin{corollary}
La transition \(\mathrm{PH}\to\mathrm{PCH}\) consiste à remplacer un espace de stabilité unique par un ensemble admissible contextuel \(\mathcal M_c\), puis à propager cette dépendance contextuelle dans l'énergie, le noyau, le résidu et la trajectoire de purification.
\end{corollary}

\subsection{Version bidirectionnelle contextuelle}

Pour rester fidèle à la purification bidirectionnelle de PH, on peut introduire une version contextuelle symétrique.

\begin{definition}[Purification bidirectionnelle contextuelle]
Soit \(u=(u_1,\dots,u_T)\) une séquence latente locale.
On définit :
\[
\mathcal F_{\alpha,c}(u)_i
=
\alpha(c)\,\mathcal F_{\alpha,c}(u)_{i-1}
+
\bigl(1-\alpha(c)\bigr)u_i,
\]
\[
\mathcal B_{\alpha,c}(u)_i
=
\alpha(c)\,\mathcal B_{\alpha,c}(u)_{i+1}
+
\bigl(1-\alpha(c)\bigr)u_i,
\]
où \(\alpha:\C\to[0,1]\) est une fonction de contrôle contextuel.
La purification contextuelle bidirectionnelle est alors
\[
\Pcal_{\alpha,c}(u)
=
\frac12\Bigl(\mathcal F_{\alpha,c}(u)+\mathcal B_{\alpha,c}(u)\Bigr).
\]
\end{definition}

\begin{remark}
Dans cette version, le contexte influence le degré de lissage, la profondeur effective de récursion, et le compromis entre compression et reconstruction.
\end{remark}

% =================================================
\section{Fonctionnelle fondamentale de stabilité contextuelle}
% =================================================

Cette section fixe la brique conceptuelle centrale de la théorie : la fonctionnelle fondamentale \(J\).
L'objectif n'est pas de poser une stabilité abstraite comme primitive, mais de la dériver à partir
d'une énergie variationnelle construite sur trois exigences observables :
(i) la fidélité à l'entrée,
(ii) la compatibilité contextuelle,
(iii) la parcimonie structurelle.

\subsection{Principe fondamental}

Dans PCH-ESM, la stabilité n'est pas une propriété primitive.
Elle résulte d'un compromis entre reconstruction, adéquation au contexte et simplicité de la représentation.

Autrement dit, la stabilité contextuelle est une quantité \emph{variationnelle} :
elle correspond à des états qui minimisent une énergie sous contrainte contextuelle.

\subsection{Contexte comme ensemble admissible}

\begin{definition}[Ensemble admissible contextuel]
À tout contexte \(c\in\C\), on associe un ensemble admissible de représentations latentes
\[
\mathcal M_c \subseteq \Z.
\]
Cet ensemble représente l'ensemble des états latents compatibles avec le contexte \(c\).

On peut le définir de manière empirique par
\[
\mathcal M_c
=
\overline{
\{\,E_\theta(x)\;:\;(x,c)\in\mathcal D\,\}
},
\]
où \(\mathcal D\) désigne l'ensemble d'entraînement ou de calibration, et la barre indique la fermeture topologique.
\end{definition}

\begin{remark}
Le contexte n'est donc pas ramené à un point central.
Il est représenté comme une région, une variété ou un ensemble admissible dans l'espace latent.
C'est cette correction qui rend la théorie robuste : un contexte n'est pas nécessairement une valeur unique,
mais une classe de compatibilité.
\end{remark}

\subsection{Fonctionnelle fondamentale}

\begin{definition}[Fonctionnelle fondamentale de stabilité]
Soient \(x\in\X\) une observation, \(c\in\C\) un contexte, et \(z\in\Z\) une représentation latente.
On définit la fonctionnelle fondamentale par
\[
J(z;x,c)
=
\lambda_{\mathrm{fid}}\,\ell_{\mathrm{fid}}(z;x,c)
+
\lambda_{\mathrm{ctx}}\,\ell_{\mathrm{ctx}}(z;c)
+
\lambda_{\mathrm{cmp}}\,\Omega(z),
\]
où :
\[
\ell_{\mathrm{fid}}(z;x,c)
=
\|D_\psi(z,c)-x\|^2
\]
mesure l'erreur de reconstruction ou de réconciliation avec l'entrée,
\[
\ell_{\mathrm{ctx}}(z;c)
=
\dist(z,\mathcal M_c)^2
=
\inf_{u\in\mathcal M_c}\|z-u\|^2
\]
mesure l'écart à l'ensemble admissible contextuel,
et \(\Omega(z)\) est une pénalisation de complexité, de variation ou de non-parcimonie.

Les coefficients \(\lambda_{\mathrm{fid}},\lambda_{\mathrm{ctx}},\lambda_{\mathrm{cmp}}>0\) pondèrent ces trois contributions.
\end{definition}

\begin{remark}
Cette définition résout la difficulté précédente sans déplacer le problème ailleurs.
Le contexte n'est pas un point caché à estimer ; il est une structure admissible dérivée des données,
et la compatibilité est mesurée par distance à cette structure.
\end{remark}

\subsection{Interprétation probabiliste équivalente}

\begin{proposition}[Forme probabiliste équivalente]
Sous une modélisation de type Gibbs, on peut associer au contexte une densité conditionnelle
\[
p_\eta(z\mid c)\propto \exp\!\left(-\frac{1}{\tau}\dist(z,\mathcal M_c)^2\right),
\]
où \(\tau>0\) est une température.
Dans cette lecture, la pénalité contextuelle peut aussi s'écrire
\[
\ell_{\mathrm{ctx}}(z;c) = -\tau \log p_\eta(z\mid c) + \text{constante}.
\]
\end{proposition}

\begin{remark}
La version géométrique et la version probabiliste sont compatibles.
La première est plus directe conceptuellement ; la seconde est utile pour l'entraînement et l'inférence.
\end{remark}

\subsection{Instabilité, stabilité et métrique dérivée}

\begin{definition}[Instabilité contextuelle fondamentale]
On définit l'instabilité contextuelle fondamentale d'un état \(z\) par
\[
\mathcal I(z;x,c)
=
J(z;x,c)
-
\inf_{u\in\Z}J(u;x,c).
\]
\end{definition}

\begin{definition}[Stabilité contextuelle fondamentale]
On définit le score de stabilité fondamentale par
\[
S_C^\star(z;x,c)
=
\frac{1}{1+\mathcal I(z;x,c)}.
\]
\end{definition}

\begin{proposition}[Bornes et interprétation]
Pour tout \(z,x,c\), on a
\[
0<S_C^\star(z;x,c)\le 1.
\]
De plus :
\begin{itemize}[leftmargin=2em]
    \item \(S_C^\star(z;x,c)=1\) si et seulement si \(z\) est un minimiseur de \(J(\cdot;x,c)\) ;
    \item \(S_C^\star(z;x,c)\) décroît lorsque l'écart au minimum contextuel augmente ;
    \item \(S_C^\star\) est donc une mesure monotone de stabilité dérivée, et non un postulat indépendant.
\end{itemize}
\end{proposition}

\begin{proof}
Par définition, \(\mathcal I(z;x,c)\ge 0\). Donc \(1+\mathcal I(z;x,c)\ge 1\), d'où
\(0<S_C^\star(z;x,c)\le 1\).
Si \(z\) est minimiseur de \(J(\cdot;x,c)\), alors \(\mathcal I(z;x,c)=0\), donc \(S_C^\star=1\).
Réciproquement, \(S_C^\star=1\) implique \(\mathcal I=0\), donc \(z\) est minimiseur.
\end{proof}

\begin{proposition}[Métrique empirique de stabilité]
En pratique, \(S_C^\star\) peut être estimé par une métrique empirique fondée sur la trajectoire de purification,
par exemple en combinant la variation de la trajectoire, la sensibilité locale du noyau et le gain d'énergie.
Cette métrique empirique est une approximation calculable de la stabilité canonique.
\end{proposition}

\subsection{Opérateur proximal canonique}

\begin{definition}[Purification contextuelle canonique]
Pour \(\lambda>0\), on définit l'opérateur de purification contextuelle par
\[
\Pcal_{\lambda}(z;x,c)
=
\argmin_{u\in\Z}
\left\{
J(u;x,c)
+
\frac{1}{2\lambda}\|u-z\|^2
\right\}.
\]
\end{definition}

\begin{remark}
Cette définition relie directement la purification au principe variationnel fondamental.
PH et PCH ne sont donc plus des mécanismes sans support énergétique explicite :
elles sont les itérations proximales de la fonctionnelle \(J\).
\end{remark}

\subsection{PH comme cas limite de la fonctionnelle fondamentale}

\begin{theorem}[PH comme cas limite de PCH]
Supposons que le contexte soit trivial au sens suivant :
\[
\mathcal M_c = \Z
\quad\text{pour tout } c\in\C,
\]
et que le décodeur soit contextuellement inactif, c'est-à-dire
\[
D_\psi(z,c)=D_\psi(z)
\quad\text{pour tout } z\in\Z,\; c\in\C.
\]
Alors la pénalité contextuelle s'annule identiquement :
\[
\ell_{\mathrm{ctx}}(z;c)=0,
\]
et la fonctionnelle fondamentale devient indépendante du contexte :
\[
J(z;x,c)=J(z;x).
\]
Par conséquent, l'opérateur \(\Pcal_{\lambda}(z;x,c)\) se réduit à l'opérateur PH classique.
\end{theorem}

\begin{proof}
Si \(\mathcal M_c=\Z\), alors
\[
\dist(z,\mathcal M_c)=0
\]
pour tout \(z\in\Z\), donc \(\ell_{\mathrm{ctx}}(z;c)=0\).
Si en outre \(D_\psi(z,c)=D_\psi(z)\), la dépendance à \(c\) disparaît complètement de \(J\).
L'opérateur proximal associé devient alors exactement celui de PH.
\end{proof}

\begin{corollary}
La transition \(\mathrm{PH}\to\mathrm{PCH}\) consiste à remplacer un espace de stabilité unique par un ensemble admissible contextuel \(\mathcal M_c\), puis à propager cette dépendance contextuelle dans l'énergie, le noyau, le résidu et la trajectoire de purification.
\end{corollary}

\begin{remark}
Cette formulation ferme la dernière faille conceptuelle :
(i) \(J\) est une énergie fondée sur des objets observables,
(ii) le contexte est un ensemble admissible, non un point arbitraire,
(iii) la stabilité est déduite comme minimum relatif,
et (iv) PH apparaît comme cas limite naturel de PCH.
\end{remark}

% =================================================
\section{PCH-ESM : espaces sémantiques modulaires et hiérarchiques}
% =================================================

Cette section formalise la décentralisation de l'espace sémantique, la spécialisation locale, la coordination globale, la gestion de la mémoire et du calcul, ainsi que l'émergence de trajectoires de pensée traçables.

\subsection{Définition générale}

\begin{definition}[Système PCH-ESM]
Un système PCH-ESM est un tuple
\[
\mathfrak M
=
\Bigl(
\E_G,
\{\E_i\}_{i=1}^n,
\{\Pi_i,\iota_i,\Pcal_i,\pi_i\}_{i=1}^n,
\Gcal,
r,
\Gamma
\Bigr),
\]
où :

\begin{itemize}[leftmargin=2em]
    \item \(\E_G\) est l'espace sémantique global ;
    \item \(\E_i\) est le \(i\)-ème espace sémantique spécialisé ;
    \item \(\Pi_i:\E_G\to\E_i\) est la projection locale ;
    \item \(\iota_i:\E_i\to\E_G\) est l'injection ou remontée vers l'espace global ;
    \item \(\Pcal_i\) est l'opérateur local de purification dans \(\E_i\) ;
    \item \(\pi_i\) est l'agent ou politique locale responsable de \(\E_i\) ;
    \item \(\Gcal\) est un graphe orienté acyclique de dépendances entre espaces ;
    \item \(r\) est une politique de routage ;
    \item \(\Gamma\) est l'opérateur de fusion globale.
\end{itemize}
\end{definition}

\subsection{Compatibilité entre espace global et espaces locaux}

\begin{assumption}[Compatibilité des projections]
Pour tout \(i\in\{1,\dots,n\}\),
\[
\Pi_i\circ \iota_i = \Id_{\E_i}.
\]
De plus, pour tout \(z\in\E_G\), l'assemblage des composantes locales reconstruit approximativement l'état global :
\[
\left\| z - \Gamma\bigl(\iota_1\Pi_1 z,\dots,\iota_n\Pi_n z\bigr)\right\|
\le \varepsilon_{\mathrm{coh}}.
\]
\end{assumption}

\begin{remark}
Cette condition exprime le fait qu'un espace global peut être décomposé en sous-espaces spécialisés sans perdre toute cohérence, à condition que la fusion et les projections soient compatibles.
\end{remark}

\subsection{Graphe hiérarchique de collaboration}

On considère un graphe orienté acyclique
\[
\Gcal=(\I,\mathcal E),
\]
dont les sommets sont les espaces spécialisés et dont les arêtes représentent des canaux de communication, de dépendance ou de transfert de contexte.

\begin{definition}[Trajectoire de pensée]
Une trajectoire de pensée est une suite finie
\[
\tau
=
\bigl((i_1,z_1,m_1),(i_2,z_2,m_2),\dots,(i_T,z_T,m_T)\bigr),
\]
où \(i_t\in\I\) est l'espace actif à l'instant \(t\), \(z_t\in\E_{i_t}\) l'état local, et \(m_t\) un message ou une mise à jour émise vers d'autres espaces.
\end{definition}

\begin{remark}
Cette définition rend la trajectoire de calcul explicitement observable. Le modèle n'est plus seulement une boîte noire qui produit une sortie ; il devient un système dont le cheminement interne peut être inspecté.
\end{remark}

\subsection{Routage et activation sparse}

\begin{definition}[Politique de routage]
La politique de routage est une application
\[
r:\E_G\times\C\to\Delta(\I),
\]
où \(\Delta(\I)\) désigne le simplexe des distributions sur \(\I\).
Pour un état global \(z\) et un contexte \(c\), la politique \(r(z,c)\) donne la probabilité ou la pondération d'activation des différents espaces.
On impose souvent une sparsité :
\[
|\supp(r(z,c))|\le k\ll n.
\]
\end{definition}

\begin{remark}
Cette sparsité formalise la gestion locale du calcul matriciel et de la mémoire : tous les espaces ne sont pas activés simultanément, seuls ceux jugés pertinents par le contexte participent à la computation.
\end{remark}

\subsection{Agents spécialisés}

\begin{definition}[Agent local]
Un agent local \(A_i\) est le triplet
\[
A_i=(\pi_i,\Pcal_i,\Theta_i),
\]
où \(\pi_i\) est la politique de décision locale, \(\Pcal_i\) l'opérateur de purification associé, et \(\Theta_i\) les paramètres internes de l'agent.
\end{definition}

Chaque agent est responsable d'un espace sémantique particulier, mais peut communiquer avec les autres agents via le graphe \(\Gcal\).

\begin{definition}[Mise à jour locale]
Si l'agent \(A_i\) est activé à l'instant \(t\), ses paramètres sont mis à jour par
\[
\Theta_i^{t+1}
=
\Theta_i^t
-
\eta_t \nabla_{\Theta_i}\mathcal L_i^t,
\]
tandis que les autres paramètres restent inchangés ou subissent une mise à jour nulle.
\end{definition}

\subsection{Traçabilité et déterminisme}

\begin{proposition}[Traçabilité]
Si la politique de routage \(r\), les opérateurs locaux \(\Pcal_i\), et l'opérateur de fusion \(\Gamma\) sont déterministes, alors la trajectoire \(\tau\) produite par le système PCH-ESM est elle-même déterministe pour une entrée et un contexte donnés.
\end{proposition}

\begin{remark}
Le déterminisme n'est donc pas automatique ; il doit être imposé. En revanche, l'architecture PCH-ESM le rend plus facile à contrôler qu'un grand espace sémantique monolithique.
\end{remark}

\subsection{Coherence locale et généralisation globale}

La spécialisation renforce la compétence locale mais fragilise la généralisation globale si aucune couche de coordination n'existe.

On distingue donc :

\begin{itemize}[leftmargin=2em]
    \item \textbf{généralisation locale} : bonne performance sur le domaine propre de l'espace spécialisé ;
    \item \textbf{généralisation globale} : bonne performance sur des tâches transversales, hors domaine initial, et sur des compositions multi-domaines.
\end{itemize}

L'architecture PCH-ESM seule ne garantit pas la généralisation globale. Celle-ci doit être reconstruite par la fonction \(G\).

% =================================================
\section{La fonction \(G\) : distillation de trajectoires et coordination inter-espaces}
% =================================================

On suppose l'existence :

\begin{itemize}[leftmargin=2em]
    \item d'un modèle généraliste \(M_G\) entraîné sur un corpus multi-disciplinaire ;
    \item d'une famille de modèles spécialisés \(\{M_i\}_{i=1}^n\), chacun entraîné sur un sous-domaine \(D_i\).
\end{itemize}

Le point clé n'est pas de distiller seulement des réponses, mais de distiller des trajectoires de coordination.

\subsection{Modèle généraliste}

Le modèle généraliste induit une politique de trajectoire
\[
\pi_G^{\mathrm{tr}}(\tau\mid x,c),
\]
où \(\tau\) est une trajectoire de pensée, de routage, de consultation d'experts et de fusion.

Le modèle généraliste apprend ainsi non seulement des connaissances, mais surtout des transitions inter-domaines.

\subsection{Modèles spécialisés}

Chaque modèle spécialisé \(M_i\) est entraîné sur un domaine \(D_i\).
Il doit être fort localement, c'est-à-dire performant sur son domaine, mais aussi apprendre à collaborer hors de son domaine.

On introduit un indicateur de domaine
\[
m_i(x)=
\begin{cases}
1 & \text{si } x\in D_i,\\
0 & \text{sinon.}
\end{cases}
\]

\subsection{Objectif de coordination}

L'objectif n'est pas que l'agent mathématique imite le comportement mathématique du généraliste sur son propre domaine, puisqu'il le maîtrise déjà localement.
L'objectif est qu'il apprenne \emph{quand et comment déléguer} ou coopérer lorsque le problème sort de son périmètre initial.

On sépare donc deux régimes :

\[
\text{si } m_i(x)=1,\quad \text{compétence locale forte};
\]
\[
\text{si } m_i(x)=0,\quad \text{coordination inter-spécialités}.
\]

\begin{definition}[Perte locale-coordination]
Pour le modèle \(M_i\), on définit
\[
\mathcal L_i
=
\mathbb E_{(x,y)\sim D_i}
\Bigl[
m_i(x)\,\mathcal L_i^{\mathrm{loc}}(x,y)
+
(1-m_i(x))\,\mathcal L_i^{\mathrm{coord}}(x)
\Bigr]
+
\lambda_i \mathcal L_i^{\mathrm{reg}}.
\]
\end{definition}

Ici :
\begin{itemize}[leftmargin=2em]
    \item \(\mathcal L_i^{\mathrm{loc}}\) mesure la compétence spécialisée ;
    \item \(\mathcal L_i^{\mathrm{coord}}\) mesure l'alignement avec les trajectoires du modèle généraliste hors domaine ;
    \item \(\mathcal L_i^{\mathrm{reg}}\) régularise le modèle ;
    \item \(\lambda_i>0\) pondère la régularisation.
\end{itemize}

\subsection{Définition de la fonction \(G\)}

\begin{definition}[Fonction de coordination \(G\)]
La fonction \(G\) est une application de distillation qui transforme les trajectoires générales en politiques locales coordonnées :
\[
G:\pi_G^{\mathrm{tr}}\longmapsto \{\pi_i^{\mathrm{coord}}\}_{i=1}^n.
\]
Plus précisément, \(G\) est obtenu par minimisation d'une divergence entre la trajectoire générale et la trajectoire produite par le système spécialisé :
\[
G^\star
=
\argmin_{G}
\;
\mathbb E_{(x,c)\sim \mathcal D}
\left[
\KL\!\left(
\pi_G^{\mathrm{tr}}(\cdot\mid x,c)
\;\middle\|\;
\pi_{G}^{\mathrm{spec}}(\cdot\mid x,c)
\right)
\right],
\]
sous contraintes de compétence locale sur chaque sous-domaine \(D_i\).
\end{definition}

\begin{remark}
La fonction \(G\) ne transfère pas la connaissance brute d'un modèle généraliste vers des spécialistes.
Elle transfère la \emph{forme des trajectoires} : quand consulter qui, dans quel ordre, avec quelle profondeur, et selon quelle hiérarchie de dépendances.
\end{remark}

\subsection{Jeu de données de coordination}

\begin{definition}[Signature de spécialité]
On note \(\I=\{1,\dots,n\}\) l'ensemble des espaces sémantiques ou spécialités.
Une signature de spécialité est une application
\[
\sigma:\X\to [0,1]^n,
\]
où \(\sigma_i(x)\) mesure le degré d'appartenance ou d'activation du prompt \(x\)
dans la spécialité \(i\).
On impose souvent une normalisation
\[
\sum_{i=1}^n \sigma_i(x)=1.
\]
\end{definition}

\begin{definition}[Trajectoire cible]
Soit \(M_G\) le modèle généraliste de référence.
On note
\[
\tau^\star(x,c)
=
\Xi(M_G,x,c)
\]
la trajectoire de référence extraite du modèle généraliste par un opérateur d'observation \(\Xi\).
Cette trajectoire encode l'ordre des espaces consultés, les passages inter-agents et les fusions intermédiaires.
\end{definition}

\begin{definition}[Trajectoire prédite]
Soit \(G_\phi\) la fonction de coordination paramétrée par \(\phi\).
On note
\[
\widehat{\tau}(x,c)=\Xi(G_\phi,x,c)
\]
la trajectoire prédite par le système coordonné.
\end{definition}

\begin{definition}[Jeu de données de coordination]
Le jeu de données utilisé pour entraîner \(G\) est
\[
\mathcal D_G
=
\bigl\{
(x,c,\sigma(x),\tau^\star(x,c),y)
\bigr\}_{k=1}^N,
\]
où :
\begin{itemize}[leftmargin=2em]
    \item \(x\) est le prompt ou l'entrée ;
    \item \(c\) est le contexte ;
    \item \(\sigma(x)\) est la signature de spécialité ;
    \item \(\tau^\star(x,c)\) est la trajectoire de référence ;
    \item \(y\) est la cible finale de la tâche.
\end{itemize}
\end{definition}

\begin{remark}
Cette construction est compatible avec l'idée suivante :
on extrait dans l'ESM la localisation des signatures de spécialité,
puis on fabrique un jeu de données prompt / signature / trajectoire,
qui révèle les chemins de coordination à apprendre par \(G\).
\end{remark}

\subsection{Fonction de perte associée à \(G\)}

\begin{definition}[Perte de tâche]
La perte de tâche standard est notée
\[
\mathcal L_{\mathrm{task}}(\hat y,y).
\]
\end{definition}

\begin{definition}[Perte de signature]
On note \(\widehat\sigma(x)\) la signature prédite par le système.
La perte de signature est définie par
\[
\mathcal L_{\mathrm{sig}}
=
\mathrm{KL}\!\bigl(\sigma(x)\,\|\,\widehat\sigma(x)\bigr),
\]
ou, si les signatures sont multi-étiquettes, par une entropie croisée binaire adaptée.
\end{definition}

\begin{definition}[Perte de trajectoire]
La perte de trajectoire est définie par
\[
\mathcal L_{\mathrm{traj}}
=
d_{\mathrm{traj}}\!\bigl(\widehat{\tau}(x,c),\tau^\star(x,c)\bigr),
\]
où \(d_{\mathrm{traj}}\) est une distance sur l'espace des trajectoires
(par exemple une divergence séquentielle, un coût d'édition, ou une somme de divergences locales le long du graphe de collaboration).
\end{definition}

\begin{definition}[Perte de cohérence de routage]
Soit \(\widehat r_t\in\Delta(\I)\) la distribution de routage prédite à l'étape \(t\).
La perte de cohérence entre routage et signature est
\[
\mathcal L_{\mathrm{route}}
=
\sum_t
\mathrm{KL}\!\bigl(r_t^\star\,\|\,\widehat r_t\bigr),
\]
où \(r_t^\star\) est le routage cible extrait du modèle généraliste ou déduit des signatures.
\end{definition}

\begin{definition}[Régularisation de graphe]
Soit \(\Gcal=(\I,\mathcal E)\) le graphe des dépendances autorisées entre espaces.
La pénalité de graphe est
\[
\mathcal L_{\mathrm{graph}}
=
\sum_t
\mathbf 1\!\left[(\widehat i_t,\widehat i_{t+1})\notin \mathcal E\right],
\]
ou une relaxation continue de cette quantité.
\end{definition}

\begin{definition}[Fonction de perte de \(G\)]
La fonction de perte totale pour \(G\) est
\[
\mathcal L_G
=
\lambda_{\mathrm{task}}\mathcal L_{\mathrm{task}}
+
\lambda_{\mathrm{sig}}\mathcal L_{\mathrm{sig}}
+
\lambda_{\mathrm{traj}}\mathcal L_{\mathrm{traj}}
+
\lambda_{\mathrm{route}}\mathcal L_{\mathrm{route}}
+
\lambda_{\mathrm{graph}}\mathcal L_{\mathrm{graph}}
+
\lambda_{\mathrm{reg}}\Omega(\phi),
\]
où \(\Omega(\phi)\) est une régularisation standard des paramètres de \(G\).
\end{definition}

\begin{remark}
Cette perte n'apprend pas seulement à produire la bonne sortie ;
elle apprend à reproduire la bonne \emph{coordination} entre espaces spécialisés.
C'est exactement ce qui permet de transférer la généralisation globale d'un modèle généraliste vers un système de spécialistes.
\end{remark}

\begin{proposition}[Interprétation de l'optimum]
Si \(\mathcal L_G\) est minimisée, alors :
\begin{enumerate}[leftmargin=2em]
    \item les signatures de spécialité sont correctement localisées ;
    \item les trajectoires internes approchent celles du modèle généraliste ;
    \item le routage respecte la structure du graphe \(\Gcal\) ;
    \item la sortie finale reste alignée avec la tâche supervisée.
\end{enumerate}
Ainsi, \(G\) apprend une politique de coordination et non une simple table de correspondance entrée-sortie.
\end{proposition}

% =================================================
\section{Multi-PH : modèle du monde et états hiérarchiques}
% =================================================

On appelle Multi-PH la version où la purification hiérarchique ne sert plus seulement à construire une représentation utile pour une tâche, mais à modéliser l'évolution d'un \emph{état du monde} latent.

\begin{definition}[État du monde hiérarchique]
On note
\[
\Omega_t = \bigl(\omega_t^{(1)},\dots,\omega_t^{(L)}\bigr)
\]
l'état du monde à l'instant \(t\), organisé en niveaux hiérarchiques.
Chaque niveau \(\omega_t^{(\ell)}\) représente une abstraction différente du monde.
\end{definition}

\begin{definition}[Dynamique PCH du monde]
La dynamique de monde est donnée par
\[
\Omega_{t+1}
=
F\!\bigl(\Omega_t,a_t,o_t\bigr),
\]
où \(a_t\) est une action, \(o_t\) une observation, et \(F\) est une application construite elle-même à partir d'une purification contextuelle hiérarchique.
\end{definition}

\begin{remark}
Dans cette perspective, prédire le prochain token n'est qu'un cas particulier d'une observation issue de l'état du monde.
La vraie variable à prédire est l'état hiérarchique \(\Omega_{t+1}\), et le token est un readout.
\end{remark}

\begin{definition}[Readout token]
Le prochain token est obtenu par une tête de lecture
\[
p_\psi(x_{t+1}\mid \Omega_{t+1})
=
D_\psi(\Omega_{t+1}).
\]
\end{definition}

% =================================================
\section{Architecture unifiée}
% =================================================

L'architecture complète peut être schématisée par
\[
x_{1:T}
\longrightarrow
E_G(x_{1:T})
\longrightarrow
Z_G
\longrightarrow
\{\E_i\}_{i=1}^n
\longrightarrow
\{\tau_i\}_{i=1}^n
\longrightarrow
\Gamma
\longrightarrow
\hat y,
\]
ou, dans la variante Multi-PH,
\[
x_{1:T}
\longrightarrow
\Omega_{1:T}
\longrightarrow
p_\psi(x_{t+1}\mid \Omega_{t+1}).
\]

\subsection{Fonction de perte globale}

On peut écrire une perte globale de la forme
\[
\mathcal L_{\mathrm{tot}}
=
\lambda_{\mathrm{task}}\mathcal L_{\mathrm{task}}
+
\lambda_{\mathrm{pch}}\mathcal L_{\mathrm{pch}}
+
\lambda_{\mathrm{coord}}\mathcal L_{\mathrm{coord}}
+
\lambda_{\mathrm{multiES}}\mathcal L_{\mathrm{multiES}}
+
\lambda_{\mathrm{world}}\mathcal L_{\mathrm{world}}
+
\lambda_{\mathrm{reg}}\Omega(\Theta,\Psi),
\]
où :
\begin{itemize}[leftmargin=2em]
    \item \(\mathcal L_{\mathrm{task}}\) est la perte de tâche ;
    \item \(\mathcal L_{\mathrm{pch}}\) force la stabilité contextuelle ;
    \item \(\mathcal L_{\mathrm{coord}}\) aligne les trajectoires entre espaces ;
    \item \(\mathcal L_{\mathrm{multiES}}\) pénalise l'incohérence entre espaces ;
    \item \(\mathcal L_{\mathrm{world}}\) entraîne le modèle du monde ;
    \item \(\Omega\) est une régularisation standard.
\end{itemize}

% =================================================
\section{Interprétation théorique}
% =================================================

La théorie PCH-ESM affirme les points suivants :

\begin{enumerate}[leftmargin=2em]
    \item La bonne représentation n'est pas seulement celle qui prédit, mais celle qui reste stable sous purification récursive dans un contexte donné.
    \item L'attention standard peut être comprise comme un cas plat de ce principe.
    \item Les poids eux-mêmes peuvent être traités comme mémoire purifiable.
    \item Un espace sémantique monolithique peut être remplacé par une hiérarchie d'espaces spécialisés.
    \item La généralisation globale ne disparaît pas nécessairement avec la spécialisation locale ; elle peut être reconstruite par une fonction de coordination \(G\) apprenant les trajectoires d'un modèle généraliste.
    \item Les trajectoires internes de calcul deviennent explicites, donc traçables et potentiellement déterministes.
    \item La version monde du formalisme (Multi-PH) remplace la prédiction de tokens comme objectif fondamental par la prédiction d'états hiérarchiques latents.
\end{enumerate}

\begin{remark}
Cette théorie ne suppose pas que la spécialisation soit gratuite : elle exige une couche de coordination pour préserver la généralisation transversale.
C'est précisément le rôle de \(G\).
\end{remark}

\section{Conclusion}

PCH-ESM fournit une formalisation unifiée de :
\begin{itemize}[leftmargin=2em]
    \item la purification contextuelle hiérarchique ;
    \item la version adaptative des paramètres ;
    \item la hiérarchisation de l'espace sémantique ;
    \item la décentralisation de l'IA en espaces spécialisés ;
    \item la coordination inter-espaces par apprentissage de trajectoires ;
    \item le modèle du monde hiérarchique.
\end{itemize}

La proposition centrale est que la structure cognitive utile n'est pas uniquement un grand espace latent global, mais un \emph{écosystème hiérarchique d'espaces sémantiques} dont la stabilité, la spécialisation et la coordination sont explicitement représentées.

\end{document}